% SPDX-License-Identifier: CC-BY-4.0
% Copyright (C) 2026 Fabian Januszewski, Abhijit (the Langlandscape team)
% Section 11: AI use. Sources: agents/rules/model-selection.md, agents/rules/opencode.md, docs/agentic-process/,
% RELEASE-NOTES.md (disclosure sentence, anonymised as approved), code/python/langlandscape/pari.py.
\paragraph{Tools.} The AI tools were Claude Code and opencode; the computational tools were PARI/GP (via \texttt{cypari2}), Python and HACS.
\begin{itemize}
\item \emph{Claude Code} (Anthropic), as a hierarchy of agents coordinated through git: a global coordinating session
  delegated all work to subagents. Models: Opus~5.5 (coordinators of stages and planning, hard Lean proofs, hard \LaTeX{} proofs,
  work-package audits), Sonnet~5.5 (routine Lean, routine \LaTeX, Python/PARI code), Haiku~5.5 (mechanical tasks: vendoring,
  logging, axiom lists, polling), Fable~5.2 (all chapter audits and Lean audits, adversarial audits, and cases where Opus failed).
  The assignment of models to task classes was set by a small controlled study during stage~01--05
  (\pathtt{agents/rules/model-selection.md}); in particular Fable was not used as an implementer. Two build machines (A and B) and a second
  team member's machine, running under a separate account, ran agent sessions too; the latter took work packages and
  \emph{cross-account blind reviews}.
\item \emph{opencode} with the models \emph{Muse Spark Open} and \emph{Mistral Large 4}: question answering and cheap first-pass
  checklists, always behind a verifier. In the planning phase Muse Spark Open was a first-pass auditor of the plans (staged plan: 7 of 70 findings fully confirmed, 22 partly, 41 rejected;
  implementation plans: 15 of 89 fully confirmed, 43 partly, 31 rejected). In stage~01 both models were tried as implementers in four
  controlled comparisons and lost each to a Claude model (Muse Spark Open failed two \LaTeX/Lean work packages, Mistral Large 4 timed out
  on a code package and passed a vendoring package, more slowly than Haiku); afterwards they were used for Q\&A only.
  \textbf{Two disclosed exceptions.} The discharge of \pathtt{Hyp.RamNrdSurj} was drafted by an external model
  (mistral-large-4-0 via opencode), run by a team member on another team member's machine, outside our agent process (where
  opencode is Q\&A-only); it was integrated into the library, independently audited (PASS WITH REMARKS) and gated from scratch
  by the agents. Likewise, the discharge of \pathtt{Hyp.BrandtStabFree} was drafted by an external model (mistral-large-4-0 via opencode) on the second team member's machine, outside our agent process; integrated into the library, independently audited (Fable, PASS) and gated by our agents.
\item \emph{PARI/GP} 2.15.4, through the Python binding \texttt{cypari2} (\pathtt{code/python/langlandscape/pari.py}) and as
  GP scripts run through \texttt{cypari2} (\pathtt{code/cas/s06-xcheck/}); Python 3.10.12. Used for class sets, Brandt
  matrices, modular-symbol spaces, newforms, Hecke algebras and congruence modules (\S\ref{sec:comp}). Every result file records
  the PARI version.
\item \emph{HACS} (C++/AVX-512/CUDA) v0.2.1, with a Jacquet--Langlands module written during the event (\S\ref{ss:hacs}).
\end{itemize}

\paragraph{What the tools did.} Within the frame designed by the two team members (\S\ref{ss:human}), the agents drafted the
chapters, implemented the Lean library, the Python code and the HACS module, ran the computations and the builds, audited
each other's output and orchestrated the work packages. The critical decisions, among them the model choice FD1 for
$\HypX$, the sign FD2, the treatment of the case $q\equiv-1\bmod p$ and the relaxed HACS protocol, were taken by the team
(list in \S\ref{ss:human}); human review is the final check.

\paragraph{Session transcripts.} Redacted transcripts of the agent sessions (Claude Code main and subagent sessions, and the
opencode sessions) are published in the repository under \pathtt{supplementary/agent-logs/}, one directory per machine, with a
manifest per machine and the redaction and post-filter tools; the README there lists the machines covered with their session counts (all four machines) and the redaction policy
(model reasoning, literature excerpts and credentials are not exported). They cover the work in this repository only. Freeze times (latest session end per machine): this machine and build machine~A 08:38\,PST on 15 November, build machine~B 22:36\,PST on 14 November, the second team member's machine 08:56\,PST on 15 November; later work is not included in the logs. A stricter post-filter then removed every tool result whose call mentions a literature-like file or tool, from the logs of all four machines (it over-removes on purpose), so that no excerpt of third-party literature is published (the development history with the unfiltered logs is not public).

\paragraph{How the output was checked.}
\begin{itemize}
\item \emph{Audits.} Every stage closed with parallel Fable audits of the chapter and of the Lean code
  (\pathtt{agents/roles/mathematics\_chapter\_audit.md}, \pathtt{lean\_audit.md}); findings were fixed before the stage was merged
  (for example, the stage-02 Lean audit found a vacuous instance of the norm hypothesis, which led to \pathtt{Hyp.RamNrdSurj};
  the stage-06 Lean audit found that \pathtt{Hyp.FreeRankTwo} is equivalent to the conclusion of its use, which is now disclosed in
  the statement). The final report of the previous integration step was audited twice (\pathtt{mathematics/audits/}).
\item \emph{Cross-account reviews.} Blind reviews by agents under a separate account on a second team member's machine, of the
  chapters of stages 01, 02, 03 and 05, of the stage-06 case $q\equiv-1\bmod p$ (second reader; it found the counterexample B2
  to a proof step), and of the previous report and the talk.
\item \emph{Lean gates.} Each stage ended with a from-scratch build on build machine A with \texttt{\#print axioms} on
  every headline declaration, a forbidden-token grep, a \texttt{sorry} count and a check that the toolchain and Mathlib are
  unchanged; from stage 02 on it was cross-checked from a fresh checkout on build machine B (the final gate of the
  \texttt{RamNrdSurj} discharge only by a quick check of the affected modules there) (\S\ref{sec:lean}).
\item \emph{Independent recomputation.} The T4 and H3 checks of stage~06 were recomputed independently in GP scripts by the
  cross-account reviewer (\res{s06-xcheck/}); HACS Lane~A reproduced the Python/PARI results of stages~01 and~05 through HACS;
  in Lane~B the native C++ quaternion core and a Python$+$\texttt{cypari2} reference agree on 69/69 orbit Smith lists.
\end{itemize}

\subsection{Human contributions and responsibility}\label{ss:human}
This subsection is the team's statement. It separates what the two team members contributed from what the AI agents did;
neither side's role is meant to be enlarged.
\begin{enumerate}[label=(\arabic*)]
\item \emph{Prior work (the seed).} The list of requirements for an integral approach is the two team members' own (prior
  work, prepared before the event): the seed document \pathtt{plan/seed/towards\_simplicial\_jl.tex}, Introduction, items
  (i)--(v) (Hecke-equivariant maps in both directions via a zig-zag of simplicial orbifolds; integrality at least over
  $\Z_{(p)}$ or $\Zp$; specialisation to classical JL over $\C$; mod $p^k$, in the limit over $p$-power levels, specialisation to
  Chenevier's $p$-adic JL; a (nearly) ordinary $\Lambda$-adic version preserving congruences), together with the caveat, after
  Ichino--Prasanna, that any integral comparison must produce, not ignore, a defect. The question (a)--(e) of
  \S\ref{ss:s00-hope} is the refined form of that list.
\item \emph{Identifying the route.} In particular, the idea of replacing the heavy arithmetic geometry of modular curves in
  the standard Ribet route (bad reduction of Shimura and modular curves) by the combinatorics of the Bruhat--Tits tree, so that
  the argument becomes computable, formalizable and easier to digest, is due to the two team members. Using Bruhat--Tits trees in
  this area is not a novel idea per se; this was preliminary work, prepared before the event (same seed document, subsection
  ``The idea of the tree route'').
\item \emph{Orchestration.} The setup for parallel mathematical, Lean and implementation development (the agent hierarchy,
  roles, rules and branches of \S\ref{sec:agentic}) was designed by the two team members. Its details were augmented by Claude
  in real time (process optimisation), together with interventions by the team members.
\item \emph{Steering and critical decisions} were made by the two team members. All times PST (UTC$-$8). Source logs: all machines on Berlin time (CET, UTC+1), report sources in UTC; converted. Examples, with the time recorded in
  the process log \pathtt{agents/status/master\_log.md} (converted to PST; process logs cover the work in this repository only) or in \pathtt{agents/status/STATE.md}
  (``Decisions''); each is a team decision:
  \begin{itemize}
  \item the Lean code is licensed under Apache-2.0 (14 November, 13:33\,PST);
  \item the frozen toolchain (Lean \texttt{v4.35.0} with a fixed Mathlib) and the policy of no \texttt{lake update}
    (team statement; the rule was committed before the agent run, 14 November, 00:08\,PST, and kept when a tool that would
    have built inside the frozen tree was requested, 14:16\,PST);
  \item FD1, the $\mu_N$-model with Picard functoriality for $\HypX$, and FD2, accepting the reading $U_q\mapsto+\Pi_q$ of
    Hsieh's sign (\S\ref{ss:hsieh}) (14 November, 17:03\,PST);
  \item for $q\equiv-1\bmod p$: first the policy ``option~2'' (apply the repairs, lift the provisional status for $M>1$ only
    after a second reader and a numerical check confirm; otherwise $M=1$ only; 17:00\,PST), then the lift of the cases (F1),
    (F1$'$), with $N(\bar\rho)>1$ left open (17:32\,PST);
  \item the relaxed HACS protocol (answers to the plan questions Q2--Q6), with ``implementation over long runs'' and
    targeted validation only, the full battery only at major milestones (16:45, 16:50 and 16:53\,PST; Q4 at 17:46\,PST);
  \item the disclosure wording for \pathtt{Hyp.FreeRankTwo} and the lemma \pathtt{Defect.freeRankTwo\_iff\_le}
    (team statement; implemented in the commit of 14 November, 17:50\,PST);
  \item to run the external draft of the \pathtt{Hyp.RamNrdSurj} discharge (opencode, outside the agent process, run by a team
    member; announced 18:22\,PST, merged after audit and gate 19:03\,PST);
  \item the narrative of the submission (Ribet's sequence central; hope $\to$ Ribet $\to$ Bruhat--Tits tree) (team statement);
  \item the structure of the submission (team statement).
  \end{itemize}
  Steering was not continuous. The team set the frame and took the critical decisions above; between interventions the agents
  worked within that frame without continuous supervision (for example, after the recorded instruction of 14 November,
  02:09\,PST, to continue through planning and implementation; process log). Decisions reserved for the team that came up
  meanwhile were taken provisionally by the orchestrating agent, marked as such, and later reviewed by the team
  (\pathtt{agents/status/STATE.md}); the examples above are those the team took.
\item \emph{Responsibility.} Responsibility for the mathematics, the Lean code, the implementation and the submission remains
  with the two team members. All final review was and is their responsibility.
\end{enumerate}
